\documentclass[11pt]{article}
\usepackage{preamble}
\setlength{\parskip}{4pt}

\begin{document}

\daybanner{AP Statistics \textperiodcentered\ Unit 2 \textperiodcentered\ Topic 2.1 \textperiodcentered\ TEACHER KEY}{Counts or Conditional Percents?}

\sectionbanner{CED TETHER}

{\small
\begin{itemize}[leftmargin=1.5em,itemsep=2pt,topsep=2pt]
  \item \textbf{VAR-1.D} Identify questions to be answered about possible relationships in data. \textbf{[Skill 1.A]}
  \item \textbf{VAR-1.D.1} Apparent patterns and associations in data may be random or not.
\end{itemize}
}
\textbf{Essential question:} When group sizes differ, which comparison reveals whether two categorical variables are related?\par
\textbf{DOK-3 skill:} 4.B \textperiodcentered\ \textbf{FRQ pattern:} {\small\texttt{counts-\allowbreak{}vs-\allowbreak{}conditional-\allowbreak{}percents-\allowbreak{}are-\allowbreak{}the-\allowbreak{}variables-\allowbreak{}related}}\par
\textbf{DOK rationale:} Part (c) weighs competing claims using the day's numerical or graphical evidence and explains a limitation or consequence; the judgment requires linked reasoning beyond a calculation.\par

\textbf{Self-paced bonus sheet.} Students take it when they choose and turn it in when they finish. Everything they need is printed on the sheet. Score part (c) only, E / P / I, by hand --- never AI-graded or auto-scored. (a) and (b) are the ladder up; use them to see where a thin (c) came from.\par

\textbf{Questions to ask}
\begin{itemize}[leftmargin=1.5em,itemsep=2pt,topsep=2pt]
  \item Which number or visible feature supports your claim?
  \item Why does the other claim go beyond that evidence?
  \item What would you tell someone who chose the other claim?
\end{itemize}

\begin{callout}{\IconWarn\ LOOK FOR}{calloutred}
\begin{itemize}[leftmargin=1.5em,itemsep=2pt,topsep=2pt]
  \item A choice with copied numbers but no explanation of how they support it.
  \item Treating a model or average as a guaranteed individual outcome.
\end{itemize}
\end{callout}

\sectionbanner{SCORING PART (c) --- E / P / I}

\textbf{Expected elements}
\begin{itemize}[leftmargin=1.5em,itemsep=2pt,topsep=2pt]
  \item \textbf{reasoning-1} (required) --- Supports Rosa using 70 percent versus 50 percent
  \item \textbf{reasoning-2} (required) --- Explains that unequal group totals make 75 versus 35 misleading
  \item \textbf{reasoning-3} (required) --- Limits the result to association, with no demonstrated cause
\end{itemize}

{\small\renewcommand{\arraystretch}{1.25}
\noindent\begin{tabular}{|p{0.5in}|p{5.9in}|}
\hline
\textbf{E} & Includes all three required reasoning elements above, with correct contextual evidence. \\ \hline
\textbf{P} & Includes at least one substantive correct reasoning link above, but omits or misstates another required element. \\ \hline
\textbf{I} & Includes no substantive correct reasoning link above; an unsupported choice or copied number alone is insufficient. \\ \hline
\end{tabular}}

\textbf{Common mistakes}
\begin{itemize}[leftmargin=1.5em,itemsep=2pt,topsep=2pt]
  \item Dividing 35 and 75 by the grand total instead of their row totals
  \item Saying 75 is a larger percentage because it is a larger count
  \item Claiming that association proves playing a sport increases sleep
\end{itemize}

\clearpage
\focusbanner{THE PROBLEM --- annotated}

Suppose a counselor records whether each of 200 students plays a school sport and whether the student usually gets at least 8 hours of sleep on a school night.\par\smallskip \textbf{Kai:} \emph{Seventy-five nonathletes but only 35 athletes get at least 8 hours, so playing a sport goes with less sleep.}\par \textbf{Rosa:} \emph{Those counts compare groups of different sizes; 35 out of 50 may tell a different story from 75 out of 150.}\par
\begin{center}
\textbf{Usual sleep on a school night}\par\smallskip
\begin{tabular}{|l|c|c|c|}
\hline
\textbf{Sport?} & \textbf{$\geq 8$ h} & \textbf{$<8$ h} & \textbf{Total} \\ \hline
\textbf{Plays} & 35 & 15 & 50 \\ \hline
\textbf{Does not} & 75 & 75 & 150 \\ \hline
\textbf{Total} & 110 & 90 & 200 \\ \hline
\end{tabular}
\end{center}
\begin{firsttakebox}{FIRST TAKE --- one sentence before you work the parts}
Which student's reasoning seems more useful for deciding whether sport participation and sleep are related? Name one table detail.\par
\answer{Not graded. Expect the 75-to-35 count comparison to attract students before they normalize for group size.}
\end{firsttakebox}

\begin{callout}{\IconBook\ COMPARE WITHIN GROUPS (keep on the page)}{calloutgreen}
A \textbf{two-way table} records combinations of two categorical variables. A raw count answers
\emph{how many}; it does not by itself answer \emph{how likely within a group}. To compare groups of
different sizes, compute a \textbf{conditional relative frequency}:
$\text{cell count}\div\text{total for that row or column}$. If the conditional distributions differ,
the variables appear associated. Association from observed data does not by itself establish cause.
\end{callout}

\dokbadge{1}~\textbf{(a)} Name the two categorical variables.\par
\answer{The variables are school-sport participation and usual school-night sleep category.}

\dokbadge{2}~\textbf{(b)} Find the percent who get at least 8 hours of sleep in each sport group. Show the two fractions.\par
\answer{Athletes: $35/50=70\%$. Nonathletes: $75/150=50\%$.}

\dokbadge{3}~\textbf{(c)} Whose reasoning is better supported? Use both percentages and explain why comparing 75 with 35 can mislead. State whether this table shows that playing a sport causes more sleep. Write 3--4 sentences.\par
\answer{Rosa's reasoning is supported: $70\%$ of athletes versus $50\%$ of nonathletes get at least 8 hours. The count 75 comes from a group of 150, three times the athlete group of 50, so the larger count does not mean a larger within-group percentage. The table shows an association for these students, but observing their choices does not show that playing a sport caused more sleep.}

\end{document}
